\documentclass[11pt]{article}

\usepackage[margin=1in]{geometry}
\usepackage[utf8]{inputenc}
\usepackage[T1]{fontenc}
\usepackage{lmodern}
\usepackage{microtype}
\usepackage{amsmath, amssymb}
\usepackage{xcolor}
\usepackage{hyperref}

\definecolor{darkblue}{rgb}{0.0, 0.0, 0.5}
\hypersetup{
  colorlinks=true,
  linkcolor=darkblue,
  citecolor=darkblue,
  urlcolor=darkblue
}

\setlength{\parskip}{0.45em}
\setlength{\parindent}{0pt}

\title{Voice versus Typing: a First-Order Model of the Speed-and-Quality\\ Compression of a Walking-Voice Agentic Writing Loop}
\author{Claude (Anthropic) \and Vivek Karmarkar}
\date{19--20 May 2026}

\begin{document}
\maketitle

\begin{abstract}
\noindent We construct a first-order mathematical model of the time-and-quality differential between two writing-loop substrates: a desk-bound typed-input loop with browser-LLM and agentic-coding assistance, and a walking-voice agentic loop in which the human dictates while walking and an AI agent handles execution. Calibrating the model against a real working day in which both substrates were carried out (the 19~May 2026 production of a multi-version short motivational paper plus a globally-available skill that packages the workflow), the model recovers an intuited equivalence---``one painful day of LLM-typing is worth roughly three hours of voice-walking work''---as a quantitative consequence of three structural facts: a fatigue-decay penalty on typing speed and correction overhead, a daily cap on continued typing because of cognitive exhaustion, and an irreducible reading-time overhead that is shared between the two substrates. With reading time admitted, the wall-clock compression is $\sim$2.72$\times$; admitting a quality-decay term sharpens the figure to a quality-adjusted compression of $\sim$4.11$\times$. The deepest finding is not the multiplier, however, but a re-identification of what the voice channel actually does: it changes the optimal division of labor between human and AI by removing the input-pain barrier, and so restores the human to an authoring role that a painful input medium structurally prevents.
\end{abstract}

\section{The question}

It is easy to claim that voice-driven agentic writing is faster than typed agentic writing. Half a day of walking and talking, with the artifacts arriving on a phone, certainly \emph{feels} faster than half a day of coffee-shop typing. But a feeling is not a compression ratio, and quality-of-life and speed are different axes that can move together, separately, or in opposition. The thought experiment that follows tries to honestly disentangle them, calibrated against the real working day in which it was conducted.

The question splits into two sub-questions. First: holding content fixed (the same paper, the same citation pool, the same number of revisions), how much faster is the voice-walking loop than the typed-desk loop, in real wall-clock hours? Second: holding wall-clock time fixed, what is the quality differential of the output produced under the two regimes? Both are needed to characterize the compression honestly, and as we will see, they couple together through a mechanism that neither sub-question reveals on its own.

\section{The freedom ladder, calibrated in days}

We assume throughout that the conceptual arc and the supporting paper-reading have already been done---both are walking-friendly activities in any era, and conflating them with the execution loop would muddle the comparison. With knowledge and ideas in hand, the question is purely about turning them into a shippable PDF.

A five-rung ladder organizes the substrate options:

L0 is the pre-LLM, desk-bound regime of roughly three years ago. The human drafts in a coffee shop, takes breaks to walk and run because the work is exhausting, manually fetches papers and BibTeX entries from Google Scholar, hand-codes TikZ for any diagrams, and grinds through revisions one painful day at a time. Honest reconstruction: nine days spread out, compressible to one week if motivated.

L1 is the browser-LLM era. The LLM helps with prose drafting and TikZ debugging, but cannot fetch or verify papers reliably. The human still hits Google Scholar by hand. Critically, the presence of the LLM raises motivation enough that the human now grinds through a full six-to-eight-hour coffee-shop session ``because the LLM is there to help,'' and the quality of life cost of that grinding is real. Honest estimate: five days, with at least one of them spent paying a substantial QoL tax.

L2 is the agentic-OS era of early 2025, with Claude Code on the laptop. The search-verify-format-compile loop collapses: the agent hits OpenAlex, Unpaywall, and Sci-Hub in parallel, assembles BibTeX, drafts prose, and debugs compilations. The human still verifies what the agent produces. Honest estimate: four days, mostly because cognitive bandwidth still caps daily output even when the typing-and-formatting toil is delegated.

L3 is Claude Code with remote-control, driven from a phone. The same agentic loop as L2, but the human is no longer desk-bound. Wall clock is roughly the same as L2 (the phone-typing penalty is offset by being able to work from anywhere), but QoL jumps sharply---the human walks twenty miles a day with a sense of constant progress. The wall-clock parity with L2 is the first sign that the speed and QoL axes can move independently.

L4 is voice over Telegram with a fully-built supporting skill ecosystem. The human dictates into Telegram voice notes; the agent transcribes, extracts, synthesizes, performs the full paper-discovery pipeline, builds LaTeX, compiles, and emails the artifact. Empirically: \emph{three hours} to a final draft, with the human walking the whole time. The day on which this thought experiment was conducted was itself an L4 day, and the empirical wall-clock matches the projection.

Tabulated, the ladder reads roughly: 10 days $\rightarrow$ 5 days $\rightarrow$ 4 days $\rightarrow$ 4 days $\rightarrow$ 3 hours. Two distinct phenomena live in this sequence. The L1$\rightarrow$L2 jump is the \emph{agentic-execution compression}, where the search-verify-format-compile loop collapses from human-serial to agent-parallel. The L3$\rightarrow$L4 jump is the \emph{cognitive-medium compression}, where voice begins to match the native cadence of broken-rambling thought in a way that typing cannot. These two compressions are different phenomena with different causes; the rest of the paper is concerned with the second.

\section{A fatigue-decay model of typed writing}

We model the typed-writing process at the granularity of a single, roughly thousand-word message. The human starts a coffee-shop session with some baseline researcher typing speed $s_0$, modest by professional standards but sustainable for original prose composition. With each successive message, two things happen. The speed decays as a function of accumulated fatigue and decision overhead:
\[
s(n) = \frac{s_0}{1 + \alpha n}, \qquad \alpha \approx 0.10.
\]
And the error-correction overhead per message grows from an initial small value:
\[
\varepsilon(n) = \varepsilon_0 + \beta n, \qquad \varepsilon_0 = 2~\text{min},~ \beta = 0.5~\text{min}.
\]
The per-message wall-clock time is then
\[
t(n) = \frac{W}{s(n)} + \varepsilon(n), \qquad W = 1000~\text{words}.
\]
We terminate the session when $t(m) > 60$~min, on the principle that taking an hour to type a thousand words means the session has become silly and should be ended.

With $s_0 = 60$~WPM (average touch-typing researcher composing original prose, not transcribing), the model produces $m = 19$ messages before the cap fires. The trajectory: $s(1) \approx 54.5$~WPM declines monotonically to $s(19) \approx 20.7$~WPM; $\varepsilon(1) = 2.5$~min grows to $\varepsilon(19) = 11.5$~min; the per-message $t(n)$ rises from 20.8~min to 59.8~min. Cumulative typing time across the 19 surviving messages: 766~min~$\approx$~12.77 hours. Total content produced: 19{,}000 words.

The decay parameter $\alpha$ dominates the outcome much more than the starting speed $s_0$. Doubling $\alpha$ from 0.05 to 0.15 cuts the total output by roughly 60\%, whereas raising $s_0$ from 60 to 80~WPM only buys a 37\% gain in surviving messages. Fatigue resistance matters more than raw speed for a sustained writing session.

\section{The daily cap and overnight reset}

The 12.77-hour figure is implausible as a single sitting---the human cannot in fact type for twelve hours straight. The honest cap is six hours of full-focus coffee-shop session, after which the human breaks and the day ends. Speed and correction overhead reset overnight; the next morning begins fresh at $s_0$. We re-run the same model with a daily cap $T_{\text{cap}} = 360$~min and a fresh-reset boundary condition.

The result is steady-state and clean: within a 6-hour day the human types 11 messages of 1000 words each, occupying 348.3~min of pure typing, with the leftover $\sim$12 minutes too short to start a 12th message. Total typed output: 11{,}000 words per day. Across two days, 22 messages and 22{,}000 words. Across an extended three-day stretched session (with mild inter-day fatigue, $s_0$ decreasing by 10\% per day): 11, then 10, then 9 messages, totaling 30 messages and 30{,}000 words.

The 11{,}000-words-per-day steady-state is consistent with the L0 estimate calibrated separately: nine days, motivated to seven, would produce 77{,}000--99{,}000 words of drafting, enough for the full polish-and-revision cycle of a 9-page short motivational paper with bibliography and one diagram. The two independent calibrations agree, which is encouraging.

\section{Reading time as irreducible overhead}

The model above tracks only the human's INPUT cost. There is also the human's READING cost---each AI-agent reply that arrives takes some minutes to read and think over. Honest estimate, calibrated by the operator: about 10 minutes per substantial reply, with shorter replies dropping under a minute and longer synthesis replies running 5-8 minutes. Ten minutes is the upper-bracket figure that accommodates skimming, deeper re-read, and the brief think-before-the-next-voice-note pause.

Reading time is the same on both sides. The voice channel doesn't make the AI's text replies arrive faster, and the reading rate of the human is not affected by whether the human's own input medium was voice or typing. So reading time is an irreducible shared overhead, and it sets a floor on the loop time that compression cannot dip below.

For a 2-day typing session producing 22 messages: typing input occupies 697 min ($\sim$76\% of the wall clock); reading the 22 replies occupies $22 \times 10 = 220$ min ($\sim$24\%); total wall clock 916.7 min~$=$~15.28 hours across the two days.

For the same 22 messages via voice (at the empirically-measured $\sim$188 WPM dictation rate): voice input occupies $22 \times 1000/188 \approx 117$~min ($\sim$35\% of the wall clock); reading the 22 replies occupies the same 220 min ($\sim$65\%); total wall clock 337 min = 5.62 hours.

The wall-clock compression ratio is
\[
\frac{T_{\text{typing}}}{T_{\text{voice}}} = \frac{916.7}{337} \approx 2.72.
\]
The 6-9$\times$ input-only ratios that look more dramatic in isolation collapse to 2.72$\times$ once reading is admitted as the shared floor it actually is. The voice channel compresses the human's INPUT phase, not the loop as a whole; as the relative weight of reading grows, the input-phase advantage shrinks proportionally.

Inverting the question---in the same 917-min wall clock as a 2-day typing session, how many voice messages can the loop sustain?---the answer is roughly 60, since each voice-message loop costs only $5.32 + 10 = 15.32$~min. That is 60{,}000 voice-words versus 22{,}000 typed words in the same wall time, which is the same 2.72$\times$ ratio expressed as content-per-time rather than time-per-content.

\section{Quality decay and the two-axis Pareto picture}

The model so far treats every message as fungible. It is not. As fatigue accumulates, the typed message's quality decays. We model this with a quality variable $Q(n)$ that drops linearly with accumulated message count:
\[
Q_{\text{typing}}(n) = \max\bigl(0,~Q_0 - \delta \cdot n\bigr), \qquad \delta = 0.02,~ Q_0 = 1.
\]
For voice, by contrast, we model $Q_{\text{voice}}(n) = Q_0$, constant. This reflects the empirical observation that walking-and-talking does not produce monotonic editorial decay in the way that hour-after-hour of typed grinding does, because the walking maintains energy and the voice channel does not impose the linearization tax that typing does.

Running the model over the realistic 3-day stretched session: 30 messages total, with per-day quality averaging 88\%, 67\%, 48\% across days one, two, and three respectively. Session-averaged typing quality: 69\%. Voice quality: 100\% throughout.

Combining the time and quality axes, we have for the same 30-message content the points
\[
\text{TYPING: } (21.74~\text{h},~69\%), \qquad \text{VOICE: } (7.66~\text{h},~100\%).
\]
The voice point Pareto-dominates the typing point: it is lower on the time axis AND higher on the quality axis. There is no axis on which typing looks better at fixed content. The quality-adjusted throughput compression is
\[
\frac{(N \cdot Q_{\text{voice}}) / T_{\text{voice}}}{(N \cdot Q_{\text{typing}}) / T_{\text{typing}}} = \frac{T_{\text{typing}}}{T_{\text{voice}}} \cdot \frac{Q_{\text{voice}}}{Q_{\text{typing}}} \approx 2.84 \times 1.45 \approx 4.11.
\]
The earlier intuited equivalence---``one painful LLM-typing day equals three voice-walking hours''---falls out of this calculation directly. A 3-day stretched typing session compresses to 7.66 voice hours; one typing day, accordingly, compresses to roughly 2.55 voice hours, which rounds to ``one painful day = three voice hours'' with no further fitting. The model recovers the operator's day-of intuition as a quantitative consequence.

\section{Painful-input versus pleasant-input regime}

The 2.72$\times$ and 4.11$\times$ figures, satisfying as they are, miss the deepest thing the voice channel actually does. The thing that became clear from the operator's own self-observation is that the input medium dictates not just speed but the OPTIMAL HUMAN-AI DIVISION OF LABOR. There are two regimes.

The painful-input regime is the typed-writing era. The input medium is exhausting, so the rational move on the human's side is to MINIMIZE human input. The human says ``here is my topic, here are some notes, here is a vague specification; agentically do the rest.'' The human keeps intent and final QA; the AI keeps ideation and execution. Output is AI-flavored, human-validated. The quality ceiling is bounded by what the human can catch in QA---which, since the human cannot audit ideation at scale, means AI-flavor errors slip through. This is exactly the failure mode that has been documented at the frontier of AI-assisted physics: the model fakes verifications, fabricates derivations, and the human in QA-only mode nearly ships them.

The pleasant-input regime is the voice-walking era. The input medium is enjoyable, so the rational move is to MAXIMIZE human input on the parts the human actually values. The human says ``here is my topic AND here is the structural arc of the argument AND here are the load-bearing citations AND here is the philosophical framing AND here is the honest caveat I want recorded.'' The human keeps intent, ideation, direction, structural taste, and QA. The AI keeps execution only---extraction, formatting, lookup, verification, LaTeX, BibTeX, compile, email. The output is human-flavored, AI-executed. The quality ceiling rises sharply because the human is directing the ideation explicitly, not handing it off behind a QA gate.

This is the cause of the quality differential the previous section modeled phenomenologically. The decay parameter $\delta$ on typed-message quality was a surface symptom; the underlying mechanism is that the painful-input medium pushes the human into a QA-only role where the ideation is being done by the wrong party. Voice puts the ideation back where it belongs.

The implication for the AI side is reciprocal and worth recording. The AI's correct role under the pleasant-input regime is NOT to one-shot the document. It is to execute against the human's direction. Where the AI adds value is in the agentic-execution compression already documented (search, verify, format, compile, email), not in the ideation. The two distinct compressions of the freedom ladder are, in retrospect, doing different jobs in the same loop: the L1$\rightarrow$L2 jump compresses what the AI does well (parallel mechanical execution), and the L3$\rightarrow$L4 jump restores what the human does well (authorial direction). Both compressions are real and both compound, but they are not duplicate measurements of the same phenomenon.

\section{Closing remarks}

The first-order model recovers the operator's intuited equivalence quantitatively and explains why the voice-walking substrate dominates the typed-desk substrate on both the time and quality axes. The deepest finding is the regime-change between painful-input and pleasant-input: voice doesn't merely compress time; it restores the human's role in ideation by removing the input-pain barrier, and the quality gain that follows is downstream of that role realignment. The 4.11$\times$ quality-adjusted compression and the 2.72$\times$ wall-clock compression are different facets of the same underlying shift.

Two honest qualifications close the analysis. First, the 3-hour voice day is contingent on a substantial pre-built skill ecosystem---Telegram voice transcription pipeline, agentic-coding scaffold, paper-discovery pipeline, LaTeX composer, email skill---without which voice would have nowhere to go. The compression ratios assume that infrastructure exists; building it from scratch was real work, performed in earlier sessions, and is not double-counted here. Second, this model applies specifically to writing-sample-scale academic production with a known arc and a known citation pool. Tasks that require precise mouse work, multi-window visual comparison, or physical-instrument operation will retain a substantial typed-desk component, and the model's compression ratios do not transfer to those.

Within the scope where it applies, however, the conclusion is sharp. The voice-walking substrate is not just faster than the typed-desk substrate. It is faster AND higher-quality, simultaneously, on both axes the operator cares about, and the cause is a role realignment that the math captures as a quality multiplier on top of a wall-clock compression.

\section*{Note on co-authorship and provenance}

This paper is the recursive use of the \texttt{/voice-writing-sample} skill on the workflow's own meta-topic. The human co-author (Vivek Karmarkar) produced the entire intellectual content of the thought experiment---the freedom ladder, the day-by-day reconstruction of each rung, the fatigue model's parameter intuitions, the daily-cap idea, the reading-time observation, the quality-decay framing, the painful-vs-pleasant regime distinction, and the closing implications---through voice messages on Telegram over the course of the same working day. The AI co-author (Claude, Anthropic Opus~4.7 in Claude Code) built the explicit mathematical model from the human's described parameters, ran the simulation, recovered the numeric ratios, and produced this LaTeX rendering. The human never sat at a desk during this thought-experiment thread and never typed a key into a keyboard for any of its content; the only non-voice human input across the broader session was the single \texttt{/goal} command used to trigger Phase~2 of an earlier paper-production run. No bibliography is included because the analysis is self-contained; no figures are included because the structure of the argument is best carried by prose and inline mathematics.

\end{document}
